%% Water Resources Research submission, typeset with AGU's own class.
%% The class and its supporting files are vendored unmodified in template/;
%% see template/PROVENANCE.md for the source and commit.
%%
%% Prose and numbers were rewritten on 2026-10-20 after the aligned-policy
%% rerun (scripts/330-335). The evidence identity of every claim is recorded in
%% results/wrr_submission_ready/major_revision_20261020/. scripts/223_wrr_build.py
%% builds this file and checks the AGU limits a desk editor checks first.

\documentclass[draft]{agujournal2025}

\usepackage{booktabs}
\usepackage{amsmath}
\usepackage{tabularx}
\usepackage{ragged2e}
\usepackage{microtype}
\ifdefined\XeTeXversion\else
  \usepackage{seqsplit}
  \let\wrrtexttt\texttt
  \renewcommand{\texttt}[1]{\wrrtexttt{\seqsplit{#1}}}
\fi
\setlength{\emergencystretch}{3em}

\newcolumntype{Y}[1]{>{\hsize=#1\hsize\raggedright\arraybackslash}X}
\newcolumntype{Z}[1]{>{\hsize=#1\hsize\raggedleft\arraybackslash}X}

\newcommand{\tabletext}{%
  \footnotesize
  \setlength{\baselineskip}{10pt}%
  \setlength{\parindent}{0pt}%
}

\begin{document}

\journalname{Water Resources Research}

\title{Hiding the peak, not missing the days: support-based screening of
incomplete river-temperature archives}

%% AUTHOR BLOCK IS A PLACEHOLDER.
%% paper_archived/author_information_required.md records why: the repository does not
%% contain verified personal information, and legal names, affiliations, ORCID
%% identifiers and funding lines must not be inferred from a version-control
%% alias. Fill metadata/submission_author_metadata.json first; the build script
%% refuses to mark this manuscript submission-ready until it is complete.
\authors{AUTHOR ONE\affil{1}}
\affiliation{1}{AFFILIATION PENDING AUTHOR CONFIRMATION}
\authoraddr{Corresponding author details pending author confirmation.}
\authorrunninghead{Ignored.}
\titlerunninghead{Ignored.}

\keypoints
  {A gap of fixed length damages the hottest 30-day mean only when it hides that window; the annual mean reacts to length}
  {At a 5\% error cap, a rule that screens hidden peak support keeps 67\% of an external archive; a gap-length rule keeps 40\%}
  {A fitted score's 7--8 point gain over that rule needs donor and residual information, not the support share; most networks need no screen}

\maketitle

\begin{abstract}
Incomplete river-temperature records are admitted by counting missing days: a
year passes if its longest gap is short enough. A gap of a given
length does not remove the same amount of the statistic computed, and in real
calendars the two vary together. We pair 5{,}114 complete station-years from 146
river networks with 53{,}721 designed outages, with 1{,}899 more that
cross gap length and position within the same station-years, and with 418 real
monitoring calendars, and score admission of the annual hottest 30-day mean. Position, not
length, orders the error: a 30-day gap over the peak window costs
0.64\,$^{\circ}$C more than the same gap outside it, while the annual mean
responds to length. The error splits exactly into the error that imputed days
inside the true window carry and a false-peak term; much of the apparent cost of
a misplaced reconstructed window is hidden share. At a 5\% cap on the error of
the admitted reconstruction, a peak-window rule keeps 67\% of a 13{,}017-case
external panel and a gap-length rule 40\%. A fitted risk score keeps 74\%, but
the same model without the peak-window share keeps 73\%: its gain over the rule
comes from donor and reconstruction diagnostics. On 418 real calendars the
unscreened error is 2.6--4.1\%, concentrated in two to four of 14 networks, and screens
refuse hot years no more often than other records. Support orders the error; a
screen's value depends on how consequential the gaps are.
\begin{plainlanguagesummary}
Rivers are warming, and managers ask how hot a river got in its hottest month.
Answering needs a year of daily temperatures; records have holes; and the usual
screening test asks how many days are missing. We show that this is the wrong
question for a summer peak. What matters is whether the missing days include the
few weeks the peak depends on. Using complete years from 146 river networks, we
place equal gaps in different positions: a gap over the peak window
spoils the hottest 30-day mean while the same gap in autumn barely moves it, and
a whole-year average behaves the other way. Screening on this hidden share keeps
two thirds of a test archive where a gap-length rule keeps two fifths. A learned risk score keeps slightly more, but the extra comes from knowing
how well neighbouring stations can fill the hole, not from the peak window. On the
holes real monitoring programmes usually leave, the peak is mostly observed, so
most networks need no screen, while the few that do fail badly. A screen refuses
records in proportion, hot years included, so it pays only where gaps are
long and badly placed.
\end{plainlanguagesummary}
\end{abstract}

\section{Introduction}

Cold-water habitat assessment, thermal-regime classification and
criterion-exceedance accounting \cite{wac2024} all reduce a year of daily water temperature to
one number --- a hottest 30-day mean, an annual mean, a seasonal amplitude
\cite{caissie2006,isaak2017} --- and it is that number, not the daily series,
that enters a management judgement. The error that matters is therefore the
error of the derived statistic. On the designed outages placed on the complete
station-years assembled here, computing the summer peak from the recorded days
alone, which is what an agency
does once a year passes a completeness test, carries a mean absolute error of
1.31\,$^{\circ}$C, and 24.9\% of station-years are then wrong by more than
1.0\,$^{\circ}$C. An error of that size turns a thermally stressful year into an
unremarkable one.

Current practice assumes that the amount of missing data decides whether the
derived number may be used, and two families of rule encode that assumption
(Table~\ref{tab:map}). Collection standards specify what must be present:
\citeA{mauger2015}, writing minimum standards for a state whose records feed
regional analyses, require a minimum sampling duration of one calendar month and
recommend a deployment period of 1 June to 31 August ``to capture the thermal
maxima''; the same standards govern which of those records enter the regional
assessment \cite{mauger2017}. Quantitative guidance constrains what may be
absent: \citeA{johnson2021} report that ``datasets with less than seven weeks of
consecutively missing data enabled estimation of all annual signal parameters
with reasonable accuracy'', expanding to ``approximately 20 weeks'' when missing
data are imputed. The gap-filling literature optimises the accuracy of the
filled series rather than the admissibility of the completed year
\cite{mohseni1998,li2017}, and treats the incomplete record as a missing-data
problem rather than as a property of the derived statistic \cite{rubin1976}, and the closest prior work fills incomplete summer
records so that a thermal peak can be computed, applies that filling uniformly,
declines no station-year, and admits July--August-only records on the
climatological position of the hottest window \cite{beaufort2022}. Large-sample
thermal models optimise daily or reach-scale prediction skill
\cite{isaak2017,corona2025,kratzert2023}, and selective-prediction methods
control error at a fixed rejection rate \cite{geifman2017,zaoui2020}, and
distribution-free risk control bounds an expectation under an exchangeability
assumption \cite{lei2018,angelopoulos2022,bai2026}. None of
these rules measures how much of the target statistic's own support a record
contains.

The assumption is plausible for some statistics and wrong for others. An annual
mean is a sum over days, so removing $k$ days removes $k$ terms and it matters
little which. A hottest 30-day mean is an extreme supported on a short window:
if the outage misses that window, the days that define the statistic are still
observed and it is recoverable however long the outage is, provided the filling
does not manufacture a higher window elsewhere; if the outage covers the window
the statistic must be inferred however short the outage is. The direction of a
completeness rule's error should therefore be opposite in the two cases, and an
average over a mixed archive would conceal the reversal. The literature is consistent with that
reversal without having isolated it: reported timing effects on annual thermal
parameters are inconsistent across regions, years and parameters
\cite{johnson2021}; reconstruction bias depends on where the target window falls
\cite{beaufort2022}; paired air--water analyses show how much of the thermal
regime is inherited from the atmosphere \cite{hare2023}; large-scale climate
modes modulate river thermal regimes \cite{sanikhani2026,you2026}; systematic, time-varying instrument bias rather than random
noise dominates long-term monitoring error \cite{jackson2025}; and
statistic-aware completeness conventions do exist in climate monitoring, but
they are defined by counts or seasons rather than by the statistic's support
\cite{zhang2011}. Irregular sampling and sampling artifacts likewise bias means
and trends and decide which model structures remain trustworthy
\cite{grey2024}, and reconstruction frameworks have been built for regulated and
data-sparse basins, including Chinese catchments and the Yangtze mainstem
\cite{huang2023,wang2025}; those studies answer estimation and reconstruction
questions, not the archive-admission question for a named annual extreme.

Testing the assumption requires separating two properties of a gap that real
records entangle. In the real gap calendars replayed here, gap length and the
share of the climatological peak window hidden correlate at $+0.71$
[0.44, 0.89] across 14 networks, and gap length correlates at $-0.90$ with the
fraction of warm-season days that remains observed; a comparison across real
gaps therefore cannot say whether a rule works because it avoids long gaps,
because it avoids badly placed ones, or because both simply track how much
warm-season record the station has. A
second obstacle is truth: in exactly the years where the admission decision
matters, the target statistic is unknown by construction, so the benefit of a
rule cannot be measured where it would be applied. A third is that evaluations
report a rule's discrimination rather than the composition of the archive it
leaves behind, and composition is what decides which catchments remain available
for analysis.

To address these limitations we assemble 5{,}114 complete river station-years
from 146 networks and nine national agencies and place 53{,}721 outages on them.
These come from two designs that compare cases across a mixed archive
(Section~2.3); the property separation itself rests on a further, smaller crossed
panel of 1{,}899 outages in which every one of 211 station-years carries all nine
combinations of three gap lengths and three positions. Admission rules are scored on held-out
river networks, on an external panel that took no part in method selection but
has since been opened, and on 418 real gap calendars transplanted from 59
incomplete station-years. Three questions organise the paper.
First, which property of a gap predicts the error of a derived thermal statistic,
and does the answer depend on which statistic is being computed? Second, how much
of that dependence is the reconstructor's doing, and which part of the error that
remains belongs to the days a gap hides and which to a false peak? Third, what
does conditioning on support change for the error among retained station-years,
which part of a fitted screen's added retention depends on support at all, and
what does refusing a record cost the hot-year judgement that the archive is used
to make? The third question decides whether the first two have practical value,
and on the evidence below they have value only in a regime we can name in
advance.

\begin{figure}[htbp]
\includegraphics[width=\textwidth]{figures/f1_support.pdf}
\caption{Hidden support, not gap length, orders the error, and the order reverses
between statistics. All panels use the 13{,}017 designed outages of the external
panel (8{,}040 single-run stratified and 4{,}977 transplanted-geometry outages)
with a day-of-year climatology reconstruction, so the panels differ only in
the statistic and in gap geometry. Cases are grouped by the \emph{true} hidden
share, an outcome-side variable, in a mixed archive in which length and true
hidden share still correlate at $+0.52$; panels (a)--(c) are therefore a
descriptive conditional comparison, and the controlled estimate is Figure~\ref{fig:reconstructor}a.
(a) Mean absolute error of the hottest 30-day
mean against the longest consecutive gap, split by the share of the true 30-day
peak window that the gap hides (grey, none; orange, partial; blue, all). Errors
are network-clustered means with 95\% intervals (2{,}000 cluster-bootstrap
draws over the 41 shared-station groups); a cell is drawn when it holds at least
25 cases. (b) The same for the annual mean. (c) Mean absolute error of five
annual statistics at the three hidden-support classes, with the same intervals;
the coldest 30-day mean moves in the opposite direction to the hottest. (d)
Distribution of hidden peak support in the designed outages (purple) and in the
418 real monitoring calendars (green). Source data:
\texttt{fig1\_source.csv}.}
\label{fig:support}
\end{figure}

\begin{table}
\caption{What each body of work optimises, decides on, and conditions on. The
rows are ordered by what they assume about the conditioning variable, not by
date.}
\label{tab:map}
\tabletext
\centering
\begin{tabularx}{\textwidth}{Y{1.15}Y{1.15}Y{0.55}Y{1.15}}
\hline
body of work & optimises & decision unit & conditioning axis \\
\hline
completeness guidance \cite{johnson2021} & tolerable amount of missing data & station-year & outage size (timing tested; effects reported inconsistent) \\
collection standards \cite{mauger2015} & comparability across a programme & deployment & season and duration \\
gap filling \cite{mohseni1998,li2017} & accuracy of the filled series & day & air temperature, donors \\
peak reconstruction \cite{beaufort2022} & spatial coverage of the thermal peak & reach & air regression applied uniformly; climatological summer window used to admit partial-season records \\
large-sample models \cite{isaak2017,corona2025,kratzert2023} & prediction skill & day or reach & meteorology, catchment attributes \\
selective prediction \cite{geifman2017,zaoui2020} & error at a fixed rejection rate & case & model confidence \\
\textbf{this work} & \textbf{error of the derived statistic among retained years} & \textbf{station-year} & \textbf{share of the statistic's support hidden} \\
\hline
\end{tabularx}
\end{table}

\section{Study system, data and design}

\subsection{The statistic and its support}

Let a station-year be a sequence $y=(y_1,\dots,y_T)$ of daily mean water
temperatures and let $M$ be the set of missing days. The annual hottest 30-day
mean is
\begin{equation}
H_{30}(y)=\max_{t\in\{30,\dots,T\}}\frac{1}{30}\sum_{s=t-29}^{t}y_s ,
\end{equation}
and its support is the set of days inside the winning window,
$S=\{s: s\in[t^\ast-29,t^\ast]\}$. For a gap $M$ we distinguish two
quantities that real records confound:
\begin{equation}
\ell=\max\{\text{length of a consecutive run in }M\},
\qquad
h=\frac{|M\cap S|}{|S|}.
\end{equation}
Here $\ell$ is what a completeness rule reads and $h$ is the share of the
statistic's own support that the record does not contain.

Two properties of this definition matter for what follows. First, $S$ is
defined by the year's own temperatures and is therefore unknown at the moment of
the admission decision. Three windows must be kept apart. The \emph{true} window
$S$ belongs to the complete reference year. The \emph{climatological} window
$S_c$ is estimated from the station's other complete years; it is the only
window a deployable rule can see, and the admission rules below use
$\hat h=|M\cap S_c|/|S_c|$. The \emph{reconstructed} window $S_r$ is the winning
window of the completed year and is an output of the reconstruction, not an
input to any decision. $S_c\neq S$ is an error in the \emph{support estimate};
$S_r\neq S$ is an error in the \emph{reconstruction}. The diagnostic panels use
the true $h$, and Section~4.3 separates the three rather than treating one as
the cause of another. Second, $h$ is
statistic-specific: for the annual mean every day is support, so $h$ degenerates
to the missing fraction, whereas for the coldest 30-day mean the support is the
winter window, and a gap that hides the summer peak is a gap that leaves the
cold support observed. The same gap therefore has three different values of $h$
and can be good for one statistic and bad for another.

\subsection{Records and panels}

Daily mean water temperature comes from nine national providers (ARSO, CHMI,
FOEN, GKD Bayern, LUBW, NVE, Rijkswaterstaat, SYKE and USGS) and, on the
external panel, from France's Hub'Eau and the UK Environment Agency. All records
are public monitoring archives reused for this analysis; no sampling was
conducted for this study. The base panel keeps networks with at least four
complete years: 146 networks, 457 stations and 5{,}114 fully observed
station-years. Two derived panels are used. The development panel holds 110
networks, 354 stations, 4{,}045 station-years and 40{,}704 designed outages; it
was used to select methods and is not a confirmation set. The external panel
holds 49 network labels, 143 distinct stations, 1{,}103 station-years and
13{,}017 designed outages. Labels that share a station are merged into
connected components before any fold is drawn, so the external panel has 41
validation groups and no target station appears in two folds; 164 station-years
appear under more than one label and stay inside one group.

The stored policy reconstruction fails the 0.5\,$^{\circ}$C tolerance on 15.3\%
of external cases and 14.9\% of development cases; those are the rates a screen
has to improve on. The real, incomplete part of the archive is 2{,}422
station-years. Those years
have no complete counterpart, so no reconstruction error can be measured on
them; they are used only to describe the gap geometry that monitoring programmes
actually leave.

\subsection{Designed outages: three designs and what each can identify}

Three designs place outages on complete station-years, and each supports a
different kind of statement. The reader should not take the 53{,}721 outages as
one factorial experiment.

\emph{Stratified outages} (33{,}749, 62.8\%) place a single contiguous run of a
target length in \{14, 30, 60, 90, 120, 180, 270\} days so that it overlaps the
\emph{climatological} peak window $S_c$ by none of it, a part (15--85\%) or all
of it. A request that no start day can satisfy is left empty: a 14-day run cannot
cover a 30-day window, and a 270-day run can hardly avoid one, so three of the
21 length-by-class cells are empty or hold fewer than 100 cases in each panel.
No scattered missing days are added; every stratified outage is one run.
\emph{Empirical outages} (19{,}972, 37.2\%) transplant the lengths and calendar
months of one to four gaps recorded by the same provider onto a complete year.
They are the closest the designed panels come to real geometry, and they carry
the entanglement of length and position with them. In the pooled panels gap
length therefore still correlates with the true hidden share at $+0.52$ (both
panels) and with the climatological overlap at $+0.45$ (development) and $+0.46$
(external), against $+0.71$ [0.44, 0.89] in the real calendars. Comparisons
across these panels, including Figures~\ref{fig:support} and
\ref{fig:reconstructor}c, are conditional comparisons in a mixed archive: they
describe what the archive contains, not what changing a gap would do.

The controlled design is the \emph{crossed panel}, generated separately on 24
development stations (it is not a subset of the 53{,}721: one third of its
outages share a realisation key with a development outage and a sixth of those
differ in geometry): 211 station-years from 24 stations in 12 networks, each
carrying all nine combinations of a 30-, 60- or 90-day run with a climatological
overlap of none, part or all, 1{,}899 outages in all, each reconstructed by five
methods fitted leave-one-year-out. Length and
climatological overlap are uncorrelated there ($0.00$ [$-0.02$, 0.03]); length
and the true hidden share remain correlated at $+0.14$ because the true window
does not coincide with $S_c$, so a longer run covers more of it within the same
class; the position contrasts below are therefore formed within station-years and
the length contrasts at a fixed true hidden share. Cells that cannot be filled are
reported as empty, not replaced by other station-years. Every identification claim
in Sections~4.1 to 4.3 rests on this panel; the admission results in
Section~4.4 rest on the stratified and empirical outages. The five reconstruction
methods see the same information on every design (Table~\ref{tab:panels}).

\subsection{Real calendars}

The 418 candidate groups of the real-calendar replay are built by transplanting
every missing day of a genuinely incomplete year onto a complete target year,
keeping the calendar position of the gap. The groups come from 59 incomplete
source station-years and fall on 67 target station-years in 14 networks, so the
count of groups overstates the independent evidence by roughly a factor of six.
This is the only setting in which the gap geometry is the one an agency actually
faces. It has no new sealed panel and it is not a confirmation set; it is a
development-period probe of what screening would do to real calendars.

\begin{table}
\caption{Panels, their units, and what each is allowed to support. ``Groups''
are connected components of networks that share a station id and are the
resampling and fold unit throughout.}
\label{tab:panels}
\tabletext
\centering
\begin{tabularx}{\textwidth}{Y{0.95}Y{0.85}Y{0.80}Y{0.55}Y{1.85}}
\hline
panel & cases & networks & groups & role \\
\hline
complete base & 5{,}114 station-years & 146 & --- & reference truth for every recovery experiment; USGS reference-network criteria follow \cite{falcone2011} \\
development & 40{,}704 designed outages & 110 & 110 & method selection and cross-fitted development evidence \\
external & 13{,}017 designed outages & 49 labels, 143 stations & 41 & cross-fitted external comparison; already opened, not a new seal \\
crossed (mechanism) & 1{,}899 outages on 211 station-years, 24 stations & 12 & 12 & the only controlled design: all nine length $\times$ position cells in every station-year; position and length contrasts, Figures~\ref{fig:reconstructor}a--b and \ref{fig:mechanism}a--c \\
real archive & 2{,}422 incomplete station-years & --- & --- & gap geometry only; no complete counterpart, so no error can be measured \\
real calendar replay & 418 groups from 59 source gaps on 67 target station-years & 14 & 14 & development-period probe of deployment on real gap calendars \\
\hline
\end{tabularx}
\end{table}

\section{Methods}

\subsection{Reconstruction and the information it may use}

Every reconstructor is fitted leave-one-year-out: the model that fills days in
year $y$ sees only other years at the same station, plus, where the method
allows it, other stations of the same network. Five reconstructors are used:
day-of-year climatology, a smoothed climatology, linear interpolation between
observed days, a ridge regression on visible donor stations, and a ridge
regression on gridded air temperature (NASA POWER \cite{nasapower2026}). A stored per-case policy selects one of
them for each station-year on validation evidence, and the reconstruction that
the admission rule accepts is that policy output, whose absolute error we write
$e$ and whose consequential failure we define as $e>0.5\,^{\circ}$C. On the
external panel the policy chooses donor regression for 69\% of cases, smoothed
climatology for 24\% and interpolation for 6\%; the climatology fallback is
where the policy fails, at 38\% against 8\% for donor regression, and it is one
reason the fitted score can separate records that share the same hidden
support. The
tolerance is fixed before any result in this paper was computed; it is roughly
half the interannual standard deviation of the hottest 30-day mean within a
station, so a failure is a year whose reconstructed rank would move.

\subsection{The admission problem}

A rule computes a score $s(x)$ from decision-time features $x$ and admits a
record when $s(x)\le\tau$. The threshold $\tau$ is selected on calibration
groups only, never on the labels of the records it will judge. We compare four
scores that are available before the decision and one reference arm:

\begin{itemize}
\item \emph{gap-length rule}: $s=\ell$, the longest consecutive missing run.
\item \emph{peak-window rule}: $s=\hat h$, the share of the climatological peak
window hidden by the gap.
\item \emph{certificate bound}: $s=b$, a calibrated upper bound on the
reconstruction error stored by a companion certificate model; it ranks records
by predicted error and carries no label from the evaluated year.
\item \emph{fitted risk score}: a quantile gradient-boosting regressor
($q=0.9$, depth 4, 350 iterations) on 22 decision-time features that include
gap geometry, visible-donor count and correlation, the reconstructor's own
training residual and the disagreement among reconstructors. It is trained on
$e$, the error of the same policy reconstruction it is used to admit, and it is
cross-fitted by group: the calibration score of each outer fold is an inner
out-of-fold prediction, so no outer test label can reach that fold's threshold.
\item \emph{reconstruct everything}: the reference arm, admitting every record
whose policy reconstruction is defined. It is what an agency does if no screen
is worth its cost, and it is the arm against which the value of screening has to
be judged.
\end{itemize}

Ties are handled once, at the score definition: a threshold is an observed value
of the raw score, and admission is $s\le\tau$, so a tie block is admitted whole
in calibration, in evaluation and at deployment. No jitter is added at any
stage. A separate fixed-retention comparison admits the safest 70\% of each
fold using a deterministic case key to break ties, which is a second deployable
rule and is reported as such.

\begin{figure}[htbp]
\includegraphics[width=\textwidth]{figures/f2_system_design.pdf}
\caption{How a station-year becomes an admissible record. Band 1: complete
station-years supply the reference truth, designed outages vary gap geometry,
five reconstructors are fitted leave-one-year-out, and the stored policy output
is the admitted object, whose error against the complete year defines a failure
at 0.5\,$^{\circ}$C. Band 2: the admission rule reads only decision-time
information, scores each record, calibrates one threshold per group fold with a
clustered bound at the 5\% cap, and admits a record whose score is at or below
the threshold; ties are admitted whole. Band 3: what the rule changes, and the
regime check that decides whether a screen is worth its cost. Schematic source:
\texttt{f2\_system\_design.drawio} (draw.io; see the Open Research statement).}
\label{fig:system}
\end{figure}

\subsection{Risk control and what it means}

The threshold is the largest candidate whose group-clustered bootstrap bound on
the calibration false-admit rate is at most 5\%, with the percentile taken at
$1-0.05/n_{\mathrm{cand}}$ over the $n_{\mathrm{cand}}$ distinct thresholds
searched. We are explicit about what this is and is not. It is an empirical,
finite-sample calibration rule for a clustered design, in the spirit of
distribution-free risk control but without its exchangeability guarantee
\cite{angelopoulos2022,bai2026}. It is not a
distribution-free guarantee for an unseen archive, and a bound that holds at
each candidate is not a bound that holds after maximising over candidates; the
Bonferroni factor is there for that reason. Section~4.5 reports what the rule
actually delivers in simulation, separating the population risk of the selected
threshold from the error rate observed in a finite test batch --- the two are not
the same event, and only the first is a property an agency can rely on.

\subsection{Uncertainty and the unit of inference}

The resampling unit is the shared-station group. All intervals are percentile
intervals from 2{,}000 group-bootstrap draws unless stated otherwise, and every
contrast is paired: the same resampled groups are used for both arms. Three fold
seeds (101, 202, 303) are pre-declared and all are reported. Shares are fractions
throughout; we write percentage points when a difference is in points. For the
hot-year analysis, the case-weighted estimate is primary and a unique
station-year weighting, in which each station-year carries total weight one
regardless of how many outage realisations it contributes, is reported as the
pre-declared sensitivity.

\subsection{Attributing the error and the gain}

Four analyses answer the attribution questions, each varying one thing.

\emph{An exact decomposition of the error.} Let $e_t=\tilde y_t-y_t$ be the
difference between the completed and the true year ($e_t=0$ on observed days),
$\bar e_W$ its mean over a 30-day window $W$, and $D_W=H_{30}(y)-\bar y_W\ge0$ the
distance of $W$ below the true maximum. Then
\begin{equation}
H_{30}(\tilde y)-H_{30}(y)=\max_W\bigl(\bar e_W-D_W\bigr)=\bar e_S+F,
\qquad F\ge 0,
\label{eq:identity}
\end{equation}
because $D_S=0$. The \emph{hiding channel} $\bar e_S$ is the error the true
window inherits from imputed days inside it and is exactly zero when none of $S$
is hidden; the \emph{false-peak channel} $F$ is how far another window, with its
own residuals added, overshoots the true one. The identity is checked on every
case. Three interventions on the recorded gap residuals, holding the truth, the
gap and the reconstructor fixed, then ask what drives each channel: permuting the
residuals within the gap (the marginal distribution and mean are kept, persistence
is destroyed; 25 permutations per case), removing the gap mean (fluctuations and
their persistence are kept), and keeping only the gap mean. These are numerical
experiments on recorded residuals, not identification of a natural process.

\emph{A ladder of information at one capacity.} The fitted score is re-learned
with the same learner, training years, nested cross-fit and calibration as in
Section~3.2, on nested subsets of its 22 features: gap and season coverage, then
donor information, then reconstruction diagnostics (training residual and
estimator disagreement), then station history, then the peak-window share
$\hat h$; the same ladder is repeated with a small logistic model. The decisive
contrast is the model that cannot read $\hat h$ against the one that can. A model
that reads the \emph{true} hidden share is not deployable and serves only to
bound how much a better support estimate could add.

\emph{Refusal against a hot-year-blind null.} A rule that ignores whether a year
is hot and keeps a share $r$ of all records refuses a share $1-r$ of the hot
years. We report $r$ minus the retention of true hot years, with a group-clustered
interval, and the loss $\rho N_{\mathrm{FN}}+N_{\mathrm{FP}}+c_{\mathrm{A}}N_{\mathrm{A}}$
of judging a hot year from the admitted reconstruction, where a refused record is
an abstention. The break-even abstention cost $c_{\mathrm{A}}^{\ast}$ is the cost
at which a screen and reconstructing everything lose the same; the screen pays
only if one abstention costs less than $c_{\mathrm{A}}^{\ast}$ wrong judgements.

\emph{An archive-level rule for whether to screen.} The share of an archive's
records in each cell of longest gap by climatological overlap is observable
before any label is seen. Weighting class-conditional failure rates of the
unscreened reconstruction, estimated on designed outages only, by that
composition predicts the archive's failure rate; the rule declares ``no screen''
when the prediction is at most the 5\% cap. Nothing in it is fitted on a real
calendar, so the real replay is an out-of-domain test of it.

\section{Results}

\subsection{The same gap has opposite consequences for different statistics}

\subsubsection{Support orders the error, and length barely moves it}

Figure~\ref{fig:support}a and b hold the gap length in view and split a mixed
archive by the true hidden support, so they describe what the archive contains.
For the hottest 30-day mean the support class sets the height of the curve and
the length moves it much less: a gap that leaves the peak window
observed costs 0.003\,$^{\circ}$C at one week and 0.04\,$^{\circ}$C at
91--180 days, while a gap that covers the window costs 0.68\,$^{\circ}$C at
15--30 days and stays near 0.9\,$^{\circ}$C from 60 days out to beyond 180. The
controlled estimate comes from the crossed panel, which holds the station-year,
the reconstructor and the length fixed and moves only the position: at a length
of 30 days the position effect is 0.64\,$^{\circ}$C [0.55, 0.74], and at 90 days
0.89\,$^{\circ}$C [0.76, 1.04] (Figure~\ref{fig:reconstructor}a), in agreement
with the 0.65\,$^{\circ}$C difference between the grouped classes at 15--30 days.
Lengthening the run from 30 to 90 days at a true hidden share of at least 0.9
adds 0.12\,$^{\circ}$C [0.06, 0.18] for the same reconstructor, one seventh of
the position effect (Figure~\ref{fig:reconstructor}b). For the annual mean the same split collapses: the three support
classes lie on top of one another and the error rises with length from
0.02\,$^{\circ}$C at one week to 0.26\,$^{\circ}$C by 180 days, with the one
class that has longer records reaching 0.41\,$^{\circ}$C.

\subsubsection{A colder statistic reverses the ordering}

The reversal is not confined to the annual mean. Figure~\ref{fig:support}c
compares five annual statistics at the three support classes. The hottest 7-day
mean behaves like the hottest 30-day mean but more steeply (0.03 against
1.13\,$^{\circ}$C). The seasonal amplitude, which is a difference of a summer and
a winter window, behaves like the hot statistics because its summer term is the
one at risk. The coldest 30-day mean moves the other way: 0.12\,$^{\circ}$C when
the gap misses the peak window against 0.07\,$^{\circ}$C when it covers it,
because a gap placed away from the summer peak is a gap that can cover the
winter minimum, and a gap over the summer peak leaves the cold support observed.
One rule, three statistics, opposite consequences: what a completeness count
cannot express is which days a statistic is made of.

Figure~\ref{fig:support}d adds the reason this matters for practice. In the
designed panel 30\% of outages hide more than half of the climatological peak
window and the median gap is 30 days. In the 418 real calendars the median gap is
7 days, the median missing fraction is 2.2\%, and 5.7\% of candidates hide more
than half of that window. The two populations are not the same population, and Section~4.4
shows what that difference does to the value of a screen.

\subsection{The reconstructor sets the height, not the shape}

The support effect is conditional on what the reconstructor can infer. Figure
\ref{fig:reconstructor}a crosses gap position with gap length at 30, 60 and 90
days for five reconstructors, within the same station-years. Hiding the climatological
peak window always costs something, but the cost ranges from
0.24\,$^{\circ}$C [0.20, 0.28] for donor-station regression to
0.64\,$^{\circ}$C [0.55, 0.74] for day-of-year climatology at 30 days. The
ordering of the methods is not stable in length: linear interpolation is the
second cheapest at 30 days (0.34\,$^{\circ}$C) and the most expensive at 90 days
(1.31\,$^{\circ}$C), while the other four keep their order.

Figure~\ref{fig:reconstructor}b asks the complementary question: with the true
hidden share held high, does length still matter? Two pre-declared cut-offs are
shown, both formed within station-years. At a share of at least 0.9
(48 station-years in 12 networks) a 90-day run costs more than a 30-day run by
0.12\,$^{\circ}$C [0.06, 0.18] for climatology, 0.14\,$^{\circ}$C [0.08, 0.20]
for smoothing, 0.12\,$^{\circ}$C [0.07, 0.18] for air regression and
0.05\,$^{\circ}$C [0.02, 0.08] for donor regression; at a share of at least 0.999
only 11 station-years in six networks qualify and the donor estimate
(0.02\,$^{\circ}$C [$-0.01$, 0.05]) is the only one whose interval includes zero. These length effects are small beside the
position effect at the same lengths, roughly 5 to 8 times smaller, but they are
distinguishable from zero. Linear interpolation is the large exception, at
0.82\,$^{\circ}$C [0.64, 1.00] and 1.12\,$^{\circ}$C [0.75, 1.45]: when the
window is entirely unobserved an interpolator must bridge the whole gap and its
error grows with the bridge, whereas a climatological or donor method substitutes
an estimate that does not depend on the length of the missing run. An earlier
reading of the same crossed panel, using a network-paired subset estimator that
needs both lengths only somewhere in the same network, put the climatology and
smoothing effects at a few hundredths of a degree with intervals that include
zero; that estimator discards the pairing within station-years and is reported
only as a sensitivity in the supporting information.

Figure~\ref{fig:reconstructor}c puts the two findings on one axis for the four
reconstructors that have a series on the external panel (that panel carries no
air-temperature series). Across the full range of hidden support every
reconstructor's error rises monotonically with the share hidden, but the curves
are at different heights and the interpolation curve bends upward at the
right-hand end. The support share is the right conditioning variable for every
method; the coefficient on it, and whether a length term survives alongside it,
are properties of the reconstructor. A claim that position dominates length as a
law would therefore be wrong, and it is not the claim we make.

\begin{figure}[htbp]
\includegraphics[width=\textwidth]{figures/f3_reconstructor.pdf}
\caption{The support effect is universal in sign and conditional in size. (a)
Difference in mean absolute error between an outage that covers the whole
climatological peak window and one that leaves it clear, within the same
station-years, at fixed length (30, 60, 90 days), for five reconstructors. (b)
Difference between a 90-day and a 30-day outage within station-years at a true
hidden share of at least 0.9 (filled; 48 station-years, 12 networks) and at least
0.999 (hollow; 11 station-years, 6 networks); only linear interpolation has a
large length effect. (c) Mean absolute error against the true hidden share of the
peak window for the four reconstructors with a series on the external designed
panel (a mixed archive; the air-temperature regression has no series there).
Panels (a) and (b) use the crossed panel of 1{,}899 outages and 12 networks with
2{,}000 network-clustered bootstrap draws; panel (c) uses the 13{,}017 external
outages with group-clustered intervals. Source data:
\texttt{fig3\_source.csv}.}
\label{fig:reconstructor}
\end{figure}

\subsection{The error splits into a hiding channel and a false-peak channel}

\subsubsection{Hiding and false peaks are separate channels with opposite signs}

Equation~\ref{eq:identity} holds to $2\times10^{-13}\,^{\circ}$C on all 9{,}449
finite cases of the crossed panel (46 donor or air fills are undefined), and
the hiding channel is exactly zero in every case whose true window is fully
observed. Figure~\ref{fig:mechanism}a splits the error by the true hidden share.
With the true window observed, the day-of-year climatology reconstruction is wrong
by 0.06\,$^{\circ}$C [0.04, 0.09] on average, all of it false peak, and in 9\% of
those cases the false peak exceeds 0.1\,$^{\circ}$C. When the window is entirely
hidden the two channels pull in opposite directions: the hidden days are filled
too cool (hiding channel $-0.75$\,$^{\circ}$C [$-0.92$, $-0.61$]) and the filling
invents a higher window elsewhere (false peak $+0.43$\,$^{\circ}$C [0.35, 0.52]),
so the mean absolute error of 0.92\,$^{\circ}$C is larger than the net signed
error. A donor regression, which carries information about the year, has the
same signs at about a fifth of the size ($-0.17$ and $+0.06$\,$^{\circ}$C). The
statement that a record whose peak window is observed is unaffected by a gap of
any length is therefore true of the hiding channel only; it holds for the error as
a whole to the extent that the filling does not manufacture a false peak, and that
depends on the reconstructor.

\subsubsection{Most of the misalignment contrast is hidden share, and persistence adds about a fifth}

Three windows are involved (Section~2.1). The climatological window $S_c$ is a
median of 12 days from the true window $S$ (90th percentile 29 days, 30 or more
days in 9.5\% of cases, identical in 4\%), so it is a noisy proxy: a gap that
covers all of $S_c$ covers all of $S$ in 44\% of cases, part of it in 52\% and
none of it in 4\%, while a gap clear of $S_c$ leaves $S$ clear in 94\%. The
reconstructed window $S_r$ is a median of 0 to 3 days from $S$ but 9 to 14 days
from $S_c$, because the filling has information about the year that $S_c$ lacks.
When $S_r$ lands away from $S$ the error is larger by 0.66, 0.63, 0.66, 0.32 and
0.49\,$^{\circ}$C for climatology, smoothing, interpolation, donor and air
regression. Most of that contrast is not misplacement. Cases whose reconstructed
window moves are overwhelmingly cases whose true window is hidden, and the hidden
share is a common cause of both. Held fixed within strata of the climatological
class and the true hidden class, the contrast falls to 0.19 [0.07, 0.26], 0.23
[0.11, 0.30], 0.13 [0.07, 0.17], 0.21 [0.14, 0.26] and 0.17\,$^{\circ}$C
[0.07, 0.25], between 35\% and 80\% lower (Figure~\ref{fig:mechanism}b). What
remains is the false-peak channel. A window switch is also rarely harmless: for
climatology 60\% of cases switch, 14\% of those within 0.1\,$^{\circ}$C of the
truth and a median 0.50\,$^{\circ}$C away, and the 7\% that switch to a window
disjoint from $S$ are a median 0.63\,$^{\circ}$C away. Within the stratum with the
true window observed a switch is by identity a false peak, so the contrast there
(0.55\,$^{\circ}$C for climatology) restates the size of the false peak and is not
independent evidence for a mechanism.

Residual persistence and the systematic offset are separable by intervention
(Figure~\ref{fig:mechanism}c). The median lag-1 autocorrelation of the residual
inside a gap is 0.84 to 0.88 for every method. Permuting the residuals within the
gap, which keeps their distribution and mean and destroys the persistence,
lowers the mean absolute error by 0.08\,$^{\circ}$C [0.05, 0.12] for climatology,
12 to 27\% depending on the reconstructor. Removing the gap mean while keeping the
fluctuations lowers it by 0.24\,$^{\circ}$C [0.19, 0.30], 49 to 59\%, and keeping
only the mean leaves three quarters of the error for climatology. Errors therefore
accumulate coherently, but the larger part of what a hidden window costs is a
systematic offset in the filled days --- the climatological and regression fills
regress towards the mean, and the year's hottest window is by construction a
maximum --- and the persistence of what is left contributes about a fifth. A
competing thermal peak does not order the error: the relation to the temperature
gap to the best disjoint window is non-monotone and flat for donor regression
(supporting information).

\subsubsection{A model that cannot read the peak-window share keeps as much as one that can}

Nothing in this section makes the support rule sufficient, and nothing yet says
what a fitted score reads that the rule does not. Figure~\ref{fig:mechanism}d adds
the information to one gradient-boosted learner in layers and, at each step,
reports the retention at the 5\% cap under the same nested cross-fit. On the
external panel the model that sees only gap geometry and season coverage keeps
63.2\% of the candidates; donor information lifts it to 68.7\%, reconstruction
diagnostics (the reconstructor's own training residual and the disagreement
among reconstructors) to 74.0\%, and station history adds nothing (73.4\%). The
peak-window share $\hat h$ is the last layer, and with it the model keeps 73.9\%.
Over three fold seeds the difference from the same model without $\hat h$ is
$+0.4$, $+0.8$ and $+0.3$ points, every interval includes zero, and on the
development panel it is $-0.5$, $+0.4$ and $+0.5$ points (Table~\ref{tab:ladder}).
The model that cannot read $\hat h$ already keeps 6.7 to 7.0 points more than the
peak-window rule (9.3 to 10.7 on development), every interval excluding zero. The
gain of the fitted score over the transparent rule therefore does not depend on
the engineered support feature.

This does not make support unimportant. A model with less to work from needs it:
added to gap and season coverage alone, $\hat h$ lifts retention by 2.4 to 5.0
points, and added to a small logistic model on six features by 6.2 to 6.8 points
(external) and 11.4 to 16.3 (development). Support is a compact and auditable
summary of what a flexible model can derive from the gap calendar, season
coverage and donor availability. The gain is also concentrated where gaps are
consequential: in the stratified outages the model without $\hat h$ keeps 61\% and
the peak-window rule 51\%, and in the transplanted real geometries 93\% and 91\%.
Finally, a better support estimate has limited room to help the fitted score. A
rule that reads the \emph{true} hidden share would keep 5.0 to 5.9 points more
than the climatological rule (8.2 to 9.3 on development), but the generic model
given the true share gains only 0.6 to 1.8 points over the one given $\hat h$
(2.8 to 4.1 on development). Locating the year's own peak better is worth
pursuing for the transparent rule; it is not where the fitted score's advantage
comes from.

\begin{table}
\caption{One learner, one training set and one calibration: retention at the 5\%
cap as information is added, mean of three fold seeds. The change columns give
the range over the seeds of the paired difference from the row above (from the
fitted score and from the peak-window rule for the last two rows), in points;
group-clustered 95\% intervals are in the supporting information. The rule rows
use no model; the italic rows read the true hidden share and are not
deployable.}
\label{tab:ladder}
\tabletext
\centering
\begin{tabularx}{\textwidth}{Y{1.9}Z{0.7}Z{0.7}Y{0.85}Y{0.85}}
\hline
model or rule & external (\%) & development (\%) & change, external (points) & change, development (points) \\
\hline
gap-length rule & 39.6 & 30.5 & --- & --- \\
peak-window rule & 66.6 & 60.4 & +27.0 & +29.9 \\
gap geometry and season coverage & 63.2 & 58.7 & $-2.0$ to $-5.3$ & $-1.2$ to $-2.1$ \\
\quad + donor information & 68.7 & 62.4 & +4.3 to +6.9 & +3.2 to +4.9 \\
\quad + reconstruction diagnostics & 74.0 & 70.2 & +4.8 to +5.7 & +6.1 to +9.0 \\
\quad + station history (generic model) & 73.4 & 70.4 & $-0.7$ to $-0.2$ & $-0.4$ to $+0.6$ \\
\quad + peak-window share $\hat h$ (fitted score) & 73.9 & 70.5 & +0.3 to +0.8 & $-0.5$ to $+0.5$ \\
certificate bound & 72.9 & 72.7 & --- & --- \\
\textit{generic model + true hidden share} & 75.0 & 73.9 & +0.6 to +1.8 & +2.8 to +4.1 \\
\textit{true-hidden-share rule} & 72.0 & 69.1 & +5.0 to +5.9 & +8.2 to +9.3 \\
\hline
\end{tabularx}
\end{table}

\begin{figure}[htbp]
\includegraphics[width=\textwidth]{figures/f4_mechanism.pdf}
\caption{Hiding, false peaks, persistence and the information a fitted score
uses. (a) Mean contribution of the hiding channel $\bar e_S$ and the false-peak
channel $F$ (Equation~\ref{eq:identity}) and mean absolute error, by true hidden
share of the peak window (0, between, 1), for the day-of-year climatology and
donor-regression reconstructions. (b) Misaligned minus aligned mean absolute
error, pooled and within strata of the climatological class and the true hidden
class, for five reconstructors. (c) Change in mean absolute error when the gap
residuals are permuted, their mean is removed, or only their mean is kept.
Panels (a)--(c) use the 1{,}899-outage crossed panel (211 station-years, 12
networks) with network-clustered 95\% intervals from 2{,}000 draws. (d) Retention
at the 5\% cap on the external panel (bars, with the range over three fold seeds)
and the development panel (circles) as information is added to one learner;
hatching marks the fitted score. Source data:
\texttt{fig4\_source.csv}.}
\label{fig:mechanism}
\end{figure}

\subsection{What it changes for the decision}

\subsubsection{Retention at the cap, and the comparison that was pre-specified}

Figure~\ref{fig:decision}a gives the retention each rule achieves at the
pre-registered 5\% cap, averaged over the three fold seeds, on the external and
development panels. On the external panel the gap-length rule keeps 39.6\% of
candidates, the peak-window rule 66.6\%, the certificate bound 72.9\% and the
fitted risk score 74.0\%. The fitted-score gain over the peak-window rule is 8.0,
7.3 and 7.0 points at seeds 101, 202 and 303 (Figure~\ref{fig:decision}d). Every
point estimate exceeds the five-point margin that was fixed before the comparison
and every interval excludes zero, but the lower limits are 5.0, 4.6 and 4.3
points, so we do not claim a gain of at least five points at 95\% confidence at
every seed. The three seeds re-split one panel, so they are repeated cross-fits
and not three confirmations. The development panel reproduces the
sign and the size: 70.4\% against 60.4\%, a gain of 9.7 to 10.3 points. Against
the certificate bound the fitted score is not reliably superior: it is $+1.6$
points [+0.2, +3.1] on the external panel at seed 101 but $+1.1$ and $+0.6$ with
intervals that include zero at the other two seeds, and $-2.6$, $-1.7$ and
$-2.5$ points on development, where the interval includes zero at all three. We
therefore report an advantage over a transparent support rule and no reliable
advantage over the calibrated bound.

At a fixed 70\% retention the ordering is the same and the scale is easier to
read: the fitted score admits 2.2\% failures on the external panel against 3.7\%
for the peak-window rule, 8.3\% for the gap-length rule and 2.8\% for the bound.
The gap-length rule is not a straw baseline that a weak comparison flatters: at
70\% retention it is the worst of the four by more than five points, which is the
cost of asking the wrong question of the archive.

\subsubsection{Refusing a record costs information, and hot years are refused like all records}

At a 20\,$^{\circ}$C hot-year threshold on the external panel, 86\% of candidates
are true hot years. Reconstructing everything refuses none of them but misjudges
1.3\% as cool. The fitted score refuses 25\% of the true hot years, the certificate
bound 26\%, the peak-window rule 33\% and the gap-length rule 60\%
(Figure~\ref{fig:decision}c). Those numbers are the retention of the rule, not a
property of hot years: a rule that ignored the temperature and kept the same share
of records would refuse the same share of hot years. The fitted score refuses 25.5\%
of all records, so hot years are refused \emph{less} often than records in general,
by 1.0 point [0.3, 1.6] at seed 101 and 0.8 to 1.0 points at the other seeds, with
the certificate bound at 1.0 and the peak-window rule at 0.1 to 0.3 points, an
interval that includes zero. On the development panel the excess refusal of the
fitted score is $-0.1$, $+0.3$ and $+0.4$ points, all intervals including zero. We
find no evidence that a screen is selective against hot years; what it costs is
the information in the records it refuses. The admitted hot-year share rises by
1.1 points under the fitted score, for a total deviation of $+0.7$ points with a
95\% interval of $-0.3$ to $+1.6$, so the archive's composition is not distorted
much either.

Whether that information is worth giving up depends on what a refused record
costs, which the error rate among admitted records does not say. Judging each
admitted year from its reconstruction and treating a refusal as an abstention,
the fitted score is wrong on 0.45\% of candidates against 1.6\% when everything
is reconstructed, at the price of abstaining on 25\%. It pays only if one
abstention costs less than 0.043 to 0.046 of a wrong hot-year judgement (0.062
to 0.063 on the development panel); if a missed hot year is five times as costly
as a false alarm the break-even is 0.15 to 0.16, and at ten times 0.28 to 0.30. A review that costs a
tenth of a wrong judgement would therefore make the fitted score worse than
reconstructing everything under equal costs, and better only if misses are
weighted heavily. The gap-length rule, which abstains on 60\%, breaks even at
0.024.

\subsubsection{Most real networks need no screen, and the gap calendar says which}

Figure~\ref{fig:decision}b opens the pooled result of the 418 real calendars. Pooled, admitting
every reconstruction puts 2.6--4.1\% of records above the tolerance, under the cap
at every point estimate, although the clustered intervals are wide (0.3 to 7.8\%
when only the target station fails). The screens keep 93--96\% of candidates and
reduce the failure rate to 0.5--1.7\%. The fitted-score advantage over the
peak-window rule shrinks to 0.2--2.4 points, and only 2 of 9 level-by-seed
intervals exclude zero. Where the peak window is usually observed, screening has
almost nothing to do: judging hot years from the unscreened reconstruction is
wrong in one of 418 groups.

The pooled rate hides where the failures are. Two networks exceed the cap when
only the target station fails (Czech Morava, 67\% on six station-years; Norwegian
basin 26, 50\% on three) and four when every donor fails; the other twelve
networks, or ten at the all-donor level, sit at 0 to 3.7\% (0 to 2.1\%). The two
failing networks are the ones whose gaps cover the peak window: the share of
records hiding more than half of the climatological window is 100\% and 50\%
there against 0 to 22\% elsewhere, and the median gap is 364 and 48 days against
2 to 13 days in ten of the other twelve. The composition rule of Section~3.5 recovers this
without a real label. Built from designed outages only, it predicts a pooled failure
rate of 3.5\% against 2.6--4.1\% observed, and it never says ``no screen'' for a
network that exceeded the cap (zero of two to four), at the price of advising a
screen for one to three of the others; it is right for 11 to 13 of 14 networks
(exact 95\% interval for the proportion 49 to 99.8\%). Among the designed
archives themselves, composition tracks the between-archive variation only
modestly (Spearman 0.35 on the 110 development networks and 0.61 on the 41
external groups; mean absolute error 5.8 against 6.6 points for the pooled rate
on development), and 90\% and 98\% of those archives exceed the cap, so the rule
advises a screen for all of them. The real evidence rests on two to four
networks that need a screen and on 3--6 station-years in each, so the rule is a
description of this replay, not a validated decision boundary: it says what to
look at --- the share of records whose gap covers the peak window --- and it has
not been shown to hold for an archive from another agency.

\begin{figure}[htbp]
\includegraphics[width=\textwidth]{figures/f5_decision.pdf}
\caption{What each rule keeps, and what the keeping costs. (a) Retention at the
5\% cap of the admitted reconstruction's error, by rule, on the external panel
(solid) and the development panel (hatched); bars are the mean of three fold
seeds and the line spans the seeds. The reference arm that admits every
reconstruction does not satisfy the cap and is omitted. ``Same model, no
support'' is the fitted score without the peak-window share. (b) Failure rate of
the unscreened reconstruction in each of the 14 real-calendar networks when only
the target station fails (filled), and the rate predicted from the network's gap
calendar by a composition rule built on designed outages only (open); labels give
the number of target station-years; the dashed line is the 5\% cap. (c) At a
20\,$^{\circ}$C hot-year threshold on the external panel, the share of true hot
years refused by each rule (solid) against the share of all records it refuses
(hatched), which is what a rule blind to temperature would refuse of the hot
years; the solid bar's line spans the three fold seeds. (d) Retention minus the
peak-window rule, with group-clustered 95\% intervals, by fold seed, for the
fitted score and for the same model without the peak-window share; the dashed
line is the pre-specified five-point margin. Source data:
\texttt{fig5\_source.csv}.}
\label{fig:decision}
\end{figure}

\subsection{The risk rule delivers the risk, and the earlier number did not}

Two properties of the selector need to be separated because an earlier version of
this analysis conflated them. The first is the population risk of the threshold
the rule selects: if the archive were infinite, would the admitted set violate
the 5\% cap? The second is the error rate observed in a finite fresh batch, which
carries its own sampling noise. We re-ran the selector on 2{,}000 synthetic
clustered archives per scenario, keeping the deployed rule identical to the one
used on the river panels, and measured both.

With the deployed simultaneous bound, the selected threshold's \emph{population}
failure rate exceeds 5\% in 0.4\% of repetitions in the baseline scenario, 0.6\%
near the cap, 0.4\% with unequal group sizes and 2.4\% with only twelve
calibration groups. The finite fresh batch exceeds 5\% far more often --- 7.8\%,
7.5\%, 4.6\% and 12.2\% in the same scenarios --- because the batch adds noise
that the rule cannot control. The 8.5\% figure quoted in the previous version of
this work was the second quantity, and reading it as the selector's error rate
was a misinterpretation; the rule is conservative, not miscalibrated. The price
is retention: the simultaneous bound keeps 18.5\% of synthetic candidates where a
pointwise bound keeps 35.9\%.

The extreme percentile the rule reads, $1-0.05/86\approx0.9994$, is numerically
delicate at 1{,}000 bootstrap draws. Re-running the selection with 10{,}000
draws changes the chosen threshold in 42 of 150 fold-by-score cells. The decision
consequences are small for the score this paper recommends --- the fitted score's
retention moves by at most 1.0 point --- and larger for the discrete rules, where
one grid step can move 3.9 points of retention. We therefore describe the
selected threshold as a calibration decision at this granularity, not as a
stable optimum.

\subsection{What the repair changed, and what it did not}

The previous version of this analysis declared the policy reconstruction as the
admitted object but calibrated the failure event on a different reconstruction.
On the external panel the two failure definitions disagree on 898 of 13{,}017
cases --- 6.9\% --- and the mismatch mattered: with the event measured on the
certificate, the fitted score's gain over the peak-window rule on the external
panel was 2.7 points; measured on the admitted reconstruction it is 7.2 points
under the same tie rule. Most of the change (about 4.5 points) comes from
measuring the right object, and about 1.3 points from admitting whole tie blocks
instead of jittering the scores. Neither repair was optional: a rule whose stated
object is the policy reconstruction cannot be calibrated on a different one, and
a rule whose calibration admits half of a tie block and whose deployment admits
all of it is two rules.

\section{Discussion}

\subsection{What is new relative to gap-length guidance}

The standing guidance is not wrong about what it was derived from. On the
estimands for which it was derived --- annual means and annual signal parameters
that integrate over the year --- we reproduce its behaviour: the error grows with
the amount of record removed and the position of the removal is second order. Our
disagreement is about the transfer of that result to a named extreme. For the
hottest 30-day mean the ordering inverts, and the inversion is not small: at the
same gap length, a badly placed 30-day outage costs about twenty times a
well-placed one, while a 180-day outage that leaves the peak window observed costs
almost nothing. A completeness rule cannot see this because the quantity that
matters --- the share of the statistic's support that is absent --- is not a
function of the number of missing days alone. Earlier work already tested timing
\cite{johnson2021} and admits partial-season records on the climatological
position of the hottest window \cite{beaufort2022}, and selective prediction
already withholds a prediction when its predicted risk is high
\cite{geifman2017,zaoui2020}; the contribution here is the attribution that
those studies did not make: a controlled position and length contrast within
station-years, an exact split of the error into a hiding and a false-peak
channel, a measurement of how much of a fitted screen's gain depends on the
support feature, and the cost of refusal against a hot-year-blind null.

This also changes what a recovery experiment has to control. A comparison across
real gaps cannot identify the support effect, because length and support are
entangled in exactly the records where the decision is made, and the pooled
designed panels are no better: length and the true hidden share still correlate
at $+0.52$ in them. The crossed panel is what makes the claim testable, and the
claim it supports is conditional: the support share is the right conditioning
variable for every reconstructor we tested, but the coefficient on it, and whether
a residual length effect survives, are properties of the reconstructor rather than
of rivers.

\subsection{Why the fitted score adds seven to ten points, and what it does not need}

The gain over the peak-window rule is real and large enough to matter for an
archive, but the ladder in Table~\ref{tab:ladder} shows where it comes from.
Donor information and the reconstructor's own diagnostics --- how well
neighbouring stations track the target, how large the training residual was, how
much the five reconstructions disagree --- supply the whole of it. The
peak-window share adds nothing once those are present. The score does read something that the
support share does not contain, and the ladder says what: information about the
reconstruction, not about the target statistic's support. The fitted score is
therefore not a more elaborate version of the support rule but a different
instrument that agrees with it, and an agency that cannot train one gives up 7 to
8 points of retention in exchange for a rule it can audit.

Neither place justifies a general claim of superiority, and the comparison with
the certificate bound says so directly. On the development panel the bound keeps
more records at the same risk; on the external panel the difference is at most
1.6 points and crosses zero at two of the three seeds. The honest reading is that a
fitted score is one way to recover the information a support rule leaves on the
table, not the only way and not demonstrably the best way. What the support rule
does better than either is explain itself: its decision is a statement about the
record, and an agency can audit it without the model. That is also the sense in
which support is the right variable: it is the best single observable, it is what
a small model needs to be competitive, and it supplies the vocabulary for the
exact decomposition above, though it is not what a flexible model needs.

\subsection{The regime boundary: when screening cannot pay}

The real-calendar result is the one that should govern adoption, and it has two
parts. The gaps that monitoring programmes actually leave are short and mostly
off-peak: a median of seven days, a median of 2.2\% of the year missing, and 5.7\%
of records hiding more than half of the peak window. In that regime the unscreened
archive sits under the cap at the point estimate, every screen removes only a
small tail, and the largest difference between the fitted score and the
transparent rule is 2.4 points. But the pooled rate averages over networks that
behave very differently: two to four of the 14 networks fail at 7 to 83\% while the
rest are essentially clean. A decision about whether to screen is therefore a
decision about records, or about the networks whose gaps cover the peak window,
and not about the pooled calendar; the share of records hiding more than half
of the climatological window separated the failing networks from the others in
this replay, but the evidence is two to four networks resting on 3 to 6
station-years each, and we do not offer it as a decision boundary.

Screening earns its cost in the other regime: long gaps, at or over the peak
window, which is where the designed panel lives, the fitted score's advantage
over the support rule is 7 to 8 points, and the support rule's advantage over
the length rule is about 27. The operational content of the paper is
therefore: compute the share of the peak window each record hides before looking
at any label; screen the records for which it is large; and judge the cost of
refusing a record against the cost of a wrong judgement before screening more.

The hot-year result makes the last point. A screen that lowers the failure rate
among admitted records while refusing 25 to 60\% of the true hot years looks
costly, but the refusal is the screen's retention, not a selection against hot
years: hot years are refused no more often than records in general. What a screen
spends is information, and whether that is worth spending depends on a cost the
data cannot supply. Under equal costs a screen pays only if a review or a re-run
of a refused record costs under about a twentieth of a wrong judgement. The two
quantities --- the error among admitted records and the abstention rate --- have
to be reported together; one without the other can always be made to look good.

\subsection{Limitations}

The external panel is already opened. Every comparison on it is a local
cross-fit: the model and the threshold for a held-out group are learned from
labelled outages on the other groups of the same panel. That answers whether an
agency that holds complete years and labelled outages for part of its network can
calibrate a screen for the rest. It does not answer whether a model and thresholds
frozen on the development archive can be carried unchanged to another agency; that
needs a sealed panel, and the three fold seeds re-split one panel and are not
three confirmations. The development panel is a sanity replication rather than an
independent test. The real-calendar replay is a development-period probe: 418
groups from 59 source gaps on 67 target station-years in 14 networks, with two to
four networks that need a screen, so its intervals are wide and its point
estimates should not be extrapolated to a national archive or to an unusually
hot or low-flow year.

The mechanism analysis rests on one controlled panel: 211 station-years from 24
stations in 12 development networks, one outage per cell, five reconstructors fitted
leave-one-year-out. The sign of the hiding channel, a cool bias in the hidden
window, is a property of fills that regress towards the mean and may differ for a
reconstructor that does not. The interventions on the residuals are numerical
experiments, not natural-process causal identification, and the permutation keeps
the residual distribution but not its dependence on the weather of the year.
The fitted score's advantage over the peak-window rule is concentrated in the
stratified outages, which are single runs, where run count and total missing days
identify the design arm; in the transplanted real geometries the two differ by two
points (93\% against 91\%).

The cell edges of the composition rule (gaps of 14, 30 and 90 days; overlap of 0
and 0.5) were fixed by us after we had seen the replay's per-network summaries
and before the rule was run; no alternative was tried, but the replay is not a
blind test of them.

The tolerance of 0.5\,$^{\circ}$C is a judgement, not a measurement. It
corresponds to about half the within-station interannual standard deviation of
the hottest 30-day mean in this panel, but an agency with a different use --- a
criterion exceedance, say --- would set a different one, and the retention
numbers scale with it. The costs in the abstention analysis are scenarios, not
measurements; we report the break-even cost rather than assume one.

The record's measurement basis is not uniform. About 74\% of development cases
are provider-declared daily means, 15\% are project-aggregated sub-daily data,
8\% have an unspecified daily definition, and 2\% are morning reference readings;
on the external panel about 94\% are declared daily means. We re-selected the
thresholds inside the declared-daily-mean subset and the ranking does not depend
on the mixed-measurement providers: the fitted score's advantage over the
peak-window rule is 8.0, 7.3 and 7.3 points on the external subset and 8.8, 8.7
and 10.4 on the development subset, close to the full-panel values, and its
comparison with the certificate bound remains interval-crossing.

Finally, the support share is defined against the year's own peak window, which
is unknown at decision time. The deployable rule uses the climatological window,
which is a median of 12 days from the true one. A rule that read the true share
would keep 5.0 to 5.9 points more than the climatological rule on the external
panel, but would add at most 1.8 points to the fitted score, so predicting the
year's window from antecedent conditions is worth building for a transparent rule
and we have not built it.

\section{Conclusions}

A completeness rule assumes that the amount of a record that is missing decides
whether the record can be used. We tested that assumption by separating what real
records entangle. The 5{,}114 complete station-years carry 53{,}721 designed
outages from two designs that compare cases across a mixed archive; a further
1{,}899 outages on 211 station-years cross gap length and position within the same
year; and 418 real monitoring calendars were replayed from 59 recorded gaps.

The share of the target statistic's own support that a gap hides orders the error
of an annual thermal extreme, whereas the amount missing orders the error of an
annual mean; the two orderings are opposite, and within station-years a 30-day
gap that covers the climatological peak window costs 0.64\,$^{\circ}$C more than
the same gap clear of it. The position effect ranges from 0.24 to 0.64\,$^{\circ}$C at 30 days
across five reconstructors; a length effect remains but is five to eight times
smaller, except for linear interpolation. The error splits exactly into a hiding channel,
a cool offset in the filled days, and a false-peak channel; most of the apparent
cost of a reconstructed window that lands elsewhere is hidden share, and residual
persistence accounts for about a fifth of the error.

For the admission decision, conditioning on support keeps 67\% of an external
panel at a 5\% error cap where a gap-length rule keeps 40\%, and a fitted risk
score keeps 74\%, 7 to 8 points more than the transparent support rule and not
reliably more than a calibrated certificate bound. The fitted score does not need
the support share to do so: the same model without it keeps 73\%, and the gain comes
from donor and reconstruction diagnostics. On the real calendars the unscreened
archive sits under the cap at the point estimate, but two to four of 14 networks
fail badly, and screens refuse hot years in proportion to all records, so what
they spend is information, not hot years. Support is the right conditioning
variable for a transparent rule and for a small model; a flexible model already
derives it. Its value is confined to archives whose gaps are consequential, and
inside that regime a transparent support rule captures most of what a complex one
does.

\section*{Open Research Statement}

All data are public monitoring records reused under their providers' terms, and
the derived tables are regenerated by the scripts listed below. The repository
holds the panel-building code, the analysis scripts, the frozen thresholds and
the figure sources; the schematic figures are draw.io sources under
\texttt{figures/wrr/drawio/}, exported to PDF with draw.io 32.5.5, and every
data figure is a Python script writing the PDF that appears here, with the table
behind each panel in
\texttt{results/wrr\_submission\_ready/major\_revision\_20261020/p0\_06\_figures/}
and
\texttt{results/wrr\_submission\_ready/major\_revision\_20261020b/r2\_06\_figures/}.
The scripts that produce every number in this manuscript are
\texttt{scripts/330} to \texttt{scripts/336} (aligned admission) and
\texttt{scripts/340} to \texttt{scripts/345} (design audit, same-capacity ladder,
three-window decomposition, hot-year null, screening-need rule and the two
figures that depend on them); \texttt{scripts/337} exports the
schematics and \texttt{scripts/338} re-checks the bibliography, with the earlier
evidence chain in \texttt{scripts/301} to \texttt{scripts/324}. Inputs, output hashes and software
versions are recorded in each run's \texttt{protocol\_fingerprint.json}. A
permanent archive with a DOI is planned at acceptance; the current repository is
a development repository and is not offered as the archival copy. Provider
licences and the processing steps that produce the daily series are documented in
\texttt{DATA\_RIGHTS.md} and the supporting information.

\section*{Acknowledgments}

We thank the national monitoring agencies listed in the supporting information
for maintaining and publishing the records this study reuses. [Funding and
author contributions pending author confirmation.]

\bibliography{references,references_additions}

\end{document}
